% FAC-EN-U0050
% French source: fac_body.tex lines 1798--1828.
% Authority: PDF physical pp. 49--50; printed pp. 244--245 (no. 52 complete).

\medskip
\phantomsection\label{fac:III:1:no:52}
\textbf{52. Cohomology of subvarieties of projective space.}
Apply \hyperref[fac:II:3:thm:4]{Theorem~4 of no.~47} to the covering
$\mathfrak{U}=\{U_i\}_{i\in I}$ defined in the \hyperref[fac:III:1:no:51]{preceding number}; this is
possible because each $U_i$ is isomorphic to $K^r$. We obtain:

\phantomsection\label{fac:III:1:prop:2}
\textsc{Proposition 2.}
\textit{If $\mathcal{F}$ is a coherent algebraic sheaf on
$X=\mathbf{P}_r(K)$, the homomorphism
\[
\sigma(\mathfrak{U}):H^n(\mathfrak{U},\mathcal{F})
  \longrightarrow H^n(X,\mathcal{F})
\]
is bijective for every $n\geq0$.}

Since $\mathfrak{U}$ consists of $r+1$ open subsets, we have (cf. \hyperref[fac:I:3:no:20:prop:2:cor]{no.~20,
corollary to Proposition~2}):

\phantomsection\label{fac:III:1:no:52:cor:1}
\textsc{Corollary.}
\[
H^n(X,\mathcal{F})=0\qquad\textit{for }n>r.
\]

This last result may be generalized as follows:

\phantomsection\label{fac:III:1:prop:3}
\textsc{Proposition 3.}
\textit{Let $V$ be an algebraic variety isomorphic to a locally closed
subvariety of a projective space $X$. Let $\mathcal{F}$ be a coherent
algebraic sheaf on $V$, and let $W$ be a subvariety of $V$ such that
$\mathcal{F}$ vanishes outside $W$. Then
\[
H^n(V,\mathcal{F})=0\qquad\text{for }n>\dim W.
\]}

In particular, on taking $W=V$, we obtain:

\phantomsection\label{fac:III:1:no:52:cor:2}
\textsc{Corollary.}
\[
H^n(V,\mathcal{F})=0\qquad\textit{for }n>\dim V.
\]

Identify $V$ with a locally closed subvariety of
$X=\mathbf{P}_r(K)$. There is an open subset $U$ of $X$ such that $V$ is
closed in $U$. We may suppose that $W$ is closed in $V$, which is plainly
permissible; then $W$ is closed in $U$. Put $F=X\setminus U$. Before proving
\hyperref[fac:III:1:prop:3]{Proposition~3}, we establish two lemmas.

\phantomsection\label{fac:III:1:lemma:1}
\textsc{Lemma 1.}
\textit{Let $k=\dim W$. There exist $k+1$ homogeneous polynomials
$P_i(t_0,\ldots,t_r)$, of degrees $>0$, which vanish on $F$ and have no
common zero on $W$.}

(By an abuse of language, a homogeneous polynomial $P$ is said to vanish at
a point $x$ of $\mathbf{P}_r(K)$ if it vanishes on $\pi^{-1}(x)$.)

We argue by induction on $k$, the case $k=-1$ being trivial. Choose one
point on each irreducible component of $W$, and let $P_1$ be a homogeneous
polynomial, of degree $>0$, which vanishes on $F$ and vanishes at none of
these points. (The existence of $P_1$ follows from the fact that $F$ is
closed, in view of the definition of the topology of $\mathbf{P}_r(K)$.)
Let $W'$ be the subvariety of $W$ formed by the points $x\in W$ such that
$P_1(x)=0$. By the construction of $P_1$, no irreducible component of $W$
is contained in $W'$, and it follows (cf. \hyperref[fac:II:1:no:36]{no.~36}) that $\dim W'<k$.
Applying the induction hypothesis to $W'$, we see that there exist $k$
homogeneous polynomials $P_2,\ldots,P_{k+1}$ which vanish on $F$ and have no
common zero on $W'$. It is clear that
$P_1,\ldots,P_{k+1}$ have the required properties.

\phantomsection\label{fac:III:1:lemma:2}
\textsc{Lemma 2.}
\textit{Let $P(t_0,\ldots,t_r)$ be a homogeneous polynomial of degree
$n>0$. The set $X_P$ of points $x\in X$ such that $P(x)\neq0$ is an affine
open subset of $X$.}

Associate with each point $y=(\mu_0,\ldots,\mu_r)\in Y$ the point of a
suitable space $K^N$ whose coordinates are all the monomials
\[
\mu_0^{m_0}\cdots\mu_r^{m_r},
\qquad m_0+\cdots+m_r=n.
\]
Passing to the quotient gives a map
\phantomsection\label{fac:III:1:no:52:map:veronese}
\[
\varphi_n:X\longrightarrow\mathbf{P}_{N-1}(K).
\]
It is classical, and moreover easy to verify, that $\varphi_n$ is a
biregular isomorphism of $X$ onto a closed subvariety of
$\mathbf{P}_{N-1}(K)$ (the ``Veronese variety''). The map $\varphi_n$
transforms the open subset $X_P$ into the locus of points of $\varphi_n(X)$
which do not lie on a certain hyperplane of $\mathbf{P}_{N-1}(K)$. Since
the complement of a hyperplane is isomorphic to an affine space, it follows
that $X_P$ is indeed isomorphic to a closed subvariety of an affine space.

We now prove \hyperref[fac:III:1:prop:3]{Proposition~3}. Extend the sheaf $\mathcal{F}$ by zero on
$U\setminus V$. We obtain a coherent algebraic sheaf on $U$, which we again
denote by $\mathcal{F}$, and we know (cf. \hyperref[fac:I:3:no:26]{no.~26}) that
\[
H^n(U,\mathcal{F})=H^n(V,\mathcal{F}).
\]
Let $P_1,\ldots,P_{k+1}$ be homogeneous polynomials satisfying the
conditions of \hyperref[fac:III:1:lemma:1]{Lemma~1}, and let $P_{k+2},\ldots,P_h$ be homogeneous
polynomials of degrees $>0$ which vanish on $W\cup F$ and have no common
zero at any point of $U\setminus W$. (To obtain such polynomials, it is
enough to take a system of homogeneous generators of the ideal defined by
$W\cup F$ in $K[t_0,\ldots,t_r]$.) For every $i$, $1\leq i\leq h$, let
$V_i$ be the set of points $x\in X$ such that $P_i(x)\neq0$. Then
$V_i\subset U$, and the preceding hypotheses show that
\phantomsection\label{fac:III:1:no:52:cover:B}
\[
\mathfrak{B}=\{V_i\}
\]
is an open covering of $U$. Moreover, \hyperref[fac:III:1:lemma:2]{Lemma~2} shows that the $V_i$ are
affine open subsets, whence
\[
H^n(\mathfrak{B},\mathcal{F})
 =H^n(U,\mathcal{F})
 =H^n(V,\mathcal{F})
\qquad\text{for every }n\geq0.
\]
On the other hand, if $n>k$ and the indices $i_0,\ldots,i_n$ are distinct,
one of these indices is $>k+1$, and $V_{i_0\cdots i_n}$ does not meet $W$.
It follows that the group of alternating cochains
\phantomsection\label{fac:III:1:no:52:formula:alternating-cochains}
\[
C^{\prime n}(\mathfrak{B},\mathcal{F})
\]
is zero if $n>k$. Therefore
$H^n(\mathfrak{B},\mathcal{F})=0$, by \hyperref[fac:I:3:no:20:prop:2]{Proposition~2 of no.~20}, which proves
\hyperref[fac:III:1:prop:3]{Proposition~3}.
